\documentclass[12pt]{article}
\usepackage[utf8]{inputenc}
\usepackage[T1]{fontenc}
\usepackage{hyperref}
\title{Vili Language Resources for K-12 Curriculum Development}
\author{Compiled by Claude Sonnet 5 (claude-sonnet-5) for Andr\'as Kornai}
\date{}
\begin{document}
\maketitle

\section{Introduction}
Vili (Civili, Kivili; ISO 639-3 vif) is a Bantu language (Zone H, Kongo group) spoken on the coast between southern Gabon and northern Cabinda. Speaker estimates vary widely, from a few thousand to about 100,000. This report rests on one web search, and the editions and details below are unverified beyond the search snippets.

\section{Primary Educational Resources}
No school textbooks or primers were found. One missionary-era catechist manual exists: \emph{Vocabulaire et premiers exercices Vili-Fran\c{c}ais sp\'ecialement destin\'es aux Cat\'echistes}, listed by a bookseller at \url{https://www.libreriagovi.com/19th-20th-century/vocabulaire-et-premiers-exercices-vili-francais-specialement-destines-aux-catechistes}.

\section{Linguistic and Reference Materials}
\begin{itemize}
\item \emph{Dictionnaire vili fran\c{c}ais} (L'Harmattan, 2008), catalogue entry: \url{https://side-ecom-prod-export.cdi.ch/en/9782296049390-dictionnaire-vili-fran%C3%A7ais}
\item Scanned PDF of a Vili--French dictionary on the Internet Archive: \url{https://ia903209.us.archive.org/19/items/DictionnaireViliFrancais/Dictionnaire_vili_fran%C3%A7ais.pdf}
\item Review article in \emph{Lexikos} placing the dictionary in the missionary lexicography tradition: \url{https://lexikos.journals.ac.za/pub/article/view/454/143}
\item Omniglot page on Vili: \url{https://omniglot.com/writing/vili.htm}
\item Vili language (Wikipedia): \url{https://en.wikipedia.org/wiki/Vili_language}
\end{itemize}

\section{Religious and Community Resources}
Omniglot credits Pierre Tchicaya de Boempire with translating the Bible into Vili from 1937. ScriptureEarth lists a Vili Bible and New Testament: \url{https://www.scriptureearth.org/00eng.php?iso=vif}. Joshua Project page: \url{https://www.joshuaproject.net/index.php/languages/vif}. Bibliographic note from the BnF: \url{https://takamtikou.bnf.fr/print/pdf/node/2997}.

\section{Digital and Audio Resources}
The Bible is reported to be available through Bible.is and YouVersion via ScriptureEarth. No other digital resources were found.

\section{Curriculum Development Context}
Vili has a Latin-script written tradition going back to the missionaries, a modern dictionary and a Bible. It is not known to be taught in schools. A curriculum would build on these texts and need new school materials.
\end{document}